\chapter{Pendahuluan}

\section{Latar Belakang}

Kawasan Malioboro adalah pusat pariwisata Kota Yogyakarta yang menampung pergerakan 
pengunjung dalam jumlah besar setiap hari. Koridor pejalan kaki sepanjang sekitar 1,3 km 
menghubungkan Titik Nol Yogyakarta, Benteng Vredeburg, Pasar Beringharjo, hingga Stasiun 
Tugu, dan menjadi ruang publik utama bagi wisatawan maupun warga lokal 
\cite{rahmah2025tingkat}. Data Pemerintah Kota Yogyakarta mencatat kunjungan ke Malioboro 
sepanjang tahun 2024 berkisar antara 118.255 pengunjung pada bulan Maret hingga 604.240 
pengunjung pada bulan Mei, dengan akumulasi sekitar 4,7 juta pengunjung dalam satu tahun 
\cite{pemkotjogja2024pengunjung}. Pada skala kota, perjalanan wisatawan nusantara ke Kota 
Yogyakarta sepanjang tahun 2025 berkisar antara 629.936 hingga 1.034.536 perjalanan per 
bulan \cite{bpsdiy2025wisnus}. Konsentrasi pengunjung memuncak pada periode tertentu. 
Saat arus mudik lebaran, volume pejalan kaki di Malioboro meningkat hampir dua kali lipat 
dibandingkan hari biasa, dengan puncak pada hari kedua dan ketiga lebaran 
\cite{lagalo2024dinamika}.

Kepadatan setinggi itu menimbulkan tekanan pada ruang dan lalu lintas kawasan. Sebagai 
pusat bisnis sekaligus destinasi wisata, Malioboro kerap mengalami kemacetan akibat 
lonjakan penggunaan kendaraan \cite{saputra2021pengelolaan}. Pengelolaan lalu lintas di 
kawasan ini bertumpu pada pengaturan administratif yang mengacu pada Peraturan Daerah 
Kota Yogyakarta Nomor 1 Tahun 2019 serta pengawasan langsung oleh petugas 
\cite{saputra2021pengelolaan}. Pemantauan kondisi kepadatan umumnya dilakukan secara 
manual melalui kamera pengawas yang diamati operator. Pendekatan ini memiliki 
keterbatasan yang melekat, yaitu cakupan pengamatan yang terbatas, kebutuhan daya yang 
tinggi, dan ketergantungan pada pengawasan operator secara terus-menerus 
\cite{khan2020crowd}. Akibatnya, informasi kepadatan yang objektif dan tersedia secara 
langsung untuk banyak titik sekaligus sulit diperoleh.

Sejumlah kota di dunia telah beralih ke sistem pemantauan berbasis kecerdasan buatan 
untuk mengatasi keterbatasan tersebut. Kota Hangzhou di Tiongkok mengoperasikan sistem 
\textit{City Brain} sejak tahun 2016, yang memproses data waktu nyata dari kamera dan 
sensor pada lampu lalu lintas untuk mengoptimalkan arus kendaraan 
\cite{caprotti2022platform}. Sistem tersebut dilaporkan meningkatkan kecepatan rata-rata 
lalu lintas sebesar 15,3 persen dan menurunkan kemacetan pada jam sibuk sebesar 9,2 
persen \cite{caprotti2022platform}. Pada tingkat teknis, \textit{computer vision} 
berbasis kamera pengawas telah terbukti mampu memantau kerumunan secara otomatis. 
Penelitian di sebuah pusat perbelanjaan di Indonesia, misalnya, menggunakan metode 
YOLOv3-Tiny untuk mendeteksi kerumunan dari rekaman kamera dan memicu peringatan ketika 
jumlah orang melampaui ambang tertentu \cite{ginting2021crowd}. Penggunaan 
\textit{computer vision} memungkinkan pengurangan ketergantungan pada pengawasan manusia 
sekaligus memperluas cakupan pemantauan \cite{khan2020crowd}.

Penerapan \textit{computer vision} di kawasan seperti Malioboro menghadapi kendala yang 
khas. Model deteksi objek pada umumnya dilatih menggunakan dataset publik berskala besar 
seperti COCO yang tidak memuat kendaraan khas Indonesia terutama di Malioboro, seperti 
andong, becak, dan bajaj. Model yang dilatih pada dataset semacam COCO maupun UA-DETRAC 
terbukti kesulitan menggeneralisasi pada kondisi lalu lintas yang padat, heterogen, dan 
tidak teratur di kota-kota berkembang \cite{sharma2025bmd}. Padahal, Malioboro memuat 
objek yang beragam, mulai dari pejalan kaki, kendaraan bermotor umum, hingga moda 
transportasi tradisional seperti andong dan becak. Kondisi ini menuntut dataset dan model 
deteksi yang dirancang khusus untuk karakteristik objek lokal yang bukan sekadar 
mengandalkan model siap pakai.

Penelitian ini membangun sistem pemantauan kepadatan pengunjung di kawasan Malioboro 
menggunakan pendekatan \textit{computer vision} dengan model YOLO. Dataset deteksi untuk 
sembilan kelas objek, yaitu andong, bajaj, becak, sepeda, bus, mobil, motor, orang, dan 
truk, dibangun dari rekaman kamera pengawas kawasan Malioboro. Beberapa varian model 
YOLOv8 dan YOLO11 dilatih dan dibandingkan untuk menentukan model dengan performa deteksi 
terbaik. Model terpilih kemudian diimplementasikan ke dalam sistem berbasis web yang 
memproses aliran video kamera pengawas secara waktu nyata dan menampilkan jumlah objek 
pada tiap kategori. Sebagai dasar penetapan ambang notifikasi kepadatan, sistem mengacu 
pada konsep tingkat pelayanan pejalan kaki (\textit{level of service}) yang diadopsi 
dalam Peraturan Menteri PUPR Nomor 03/PRT/M/2014 
\cite{fruin1971designing, pupr2014pedoman}. Berdasarkan uraian tersebut, beberapa 
permasalahan dirumuskan pada subbab berikut.

\section{Rumusan Masalah}

Berdasarkan latar belakang yang telah diuraikan, yaitu kebutuhan pemantauan kepadatan 
otomatis di kawasan Malioboro yang masih dikelola secara manual, permasalahan dalam 
penelitian ini dirumuskan sebagai berikut.

\begin{enumerate}
    \item Dataset deteksi objek untuk sembilan kelas khas kawasan Malioboro, yaitu andong, bajaj, becak, sepeda, bus, mobil, motor, orang, dan truk, belum tersedia karena sebagian kelas tersebut tidak terdapat pada dataset publik berskala besar seperti COCO.
    \item Varian model YOLOv8 dan YOLO11 yang memberikan performa deteksi terbaik untuk kesembilan kelas tersebut belum diketahui, termasuk pengaruh ketidakseimbangan jumlah \textit{instance} antar kelas terhadap performa deteksi pada tiap kelas.
    \item Model deteksi terbaik belum diimplementasikan ke dalam sistem pemantauan kepadatan berbasis web yang memproses aliran video kamera pengawas secara waktu nyata.
\end{enumerate}

\section{Tujuan Penelitian}

Tujuan penelitian ini dirumuskan untuk menjawab setiap permasalahan yang telah dipaparkan, 
dengan keterkaitan satu-satu antara tiap tujuan dan tiap rumusan masalah. Ketiga tujuan 
tersusun mengikuti alur penelitian, mulai dari pembentukan dataset sebagai fondasi, 
dilanjutkan pemilihan dan optimasi model deteksi melalui perbandingan dan penyetelan, 
hingga penerapan model terpilih ke dalam sistem yang berfungsi.

\begin{enumerate}
    \item Membangun dataset deteksi sembilan kelas objek kawasan Malioboro (andong, bajaj, becak, sepeda, bus, mobil, motor, orang, dan truk) dari rekaman kamera pengawas melalui anotasi semi-otomatis berbantuan model \textit{pretrained} dengan koreksi manual.
    \item Membandingkan performa varian model YOLOv8 dan YOLO11 dalam mendeteksi kesembilan kelas menggunakan metrik mAP50, mAP50-95, dan F1, menganalisis pengaruh ketidakseimbangan kelas terhadap performa deteksi tiap kelas, lalu memilih model terbaik dan mengoptimasi performanya melalui penyetelan \textit{hyperparameter}.
    \item Mengimplementasikan model deteksi terbaik ke dalam sistem pemantauan kepadatan berbasis web yang memproses aliran video kamera pengawas secara waktu nyata dan mengukur performa inferensinya.
\end{enumerate}

\section{Batasan Penelitian}

Penelitian ini diberi batasan agar lingkupnya tetap terarah dan dapat diselesaikan 
dengan sumber daya yang tersedia. Batasan mencakup lokasi dan sumber data, kelas objek 
yang dideteksi, varian model yang dibandingkan, serta fungsi sistem yang dikembangkan. 
Penetapan batasan ini sekaligus menjaga agar perbandingan model dilakukan pada kondisi 
yang terkontrol. Ruang lingkup selengkapnya diuraikan sebagai berikut.

\begin{enumerate}
    \item Lokasi penelitian dibatasi pada kawasan wisata Malioboro, Kota Yogyakarta, dengan sumber data berupa aliran video dari 22 titik kamera pengawas (\textit{closed-circuit television}, CCTV) publik di kawasan tersebut.
    \item Objek yang dideteksi dibatasi pada sembilan kelas, yaitu andong, bajaj, becak, sepeda, bus, mobil, motor, orang, dan truk.
    \item Model deteksi yang dilatih dan dibandingkan dibatasi pada varian YOLOv8 dan YOLO11.
    \item Dataset yang digunakan adalah data primer yang dibangun dari rekaman kamera pengawas kawasan Malioboro melalui anotasi semi-otomatis dengan koreksi manual. Rincian proses anotasi dijelaskan pada Bab III.
    \item Sistem hanya melakukan pemantauan (\textit{monitoring}) kepadatan secara waktu nyata dan tidak melakukan prediksi kepadatan pada waktu mendatang.
    \item Penghitungan objek dilakukan per \textit{frame}. Sistem tidak melakukan pelacakan (\textit{tracking}) objek antar-\textit{frame}.
    \item Penetapan ambang notifikasi berbasis tingkat pelayanan (\textit{level of service}) untuk kepadatan orang hanya diterapkan pada dua titik kamera yang dikalibrasi sebagai contoh. Titik kamera lainnya serta penetapan ambang untuk kemacetan kendaraan berada di luar lingkup dan menjadi arah pengembangan lanjutan.
    \item Penelitian berfokus pada perbandingan model deteksi dan implementasi fungsional sistem. Perancangan berbasis \textit{User-Centered Design} (UCD) dan pengujian \textit{usability} tidak termasuk dalam lingkup penelitian.
\end{enumerate}

\section{Manfaat Penelitian}

Hasil penelitian ini diharapkan bermanfaat pada dua sisi yang saling melengkapi, yaitu 
pengembangan keilmuan dan penerapan di lapangan. Sisi akademis berkaitan dengan 
ketersediaan data dan temuan yang dapat dijadikan rujukan penelitian berikutnya, 
sedangkan sisi praktis berkaitan dengan kegunaan sistem bagi pengelola kawasan Malioboro. 
Uraian masing-masing manfaat adalah sebagai berikut.

\noindent Manfaat akademis:

\begin{enumerate}
    \item Penelitian ini menghasilkan dataset deteksi objek untuk sembilan kelas khas kawasan Malioboro, termasuk andong, becak, dan bajaj yang tidak terdapat pada dataset publik seperti COCO, sehingga dapat menjadi sumber data bagi penelitian deteksi objek lokal selanjutnya.
    \item Hasil perbandingan performa varian model YOLOv8 dan YOLO11 pada dataset objek lokal dengan distribusi kelas yang bervariasi dapat menjadi rujukan dalam pemilihan model deteksi objek untuk kasus serupa.
\end{enumerate}

\noindent Manfaat praktis:

\begin{enumerate}
    \item Sistem pemantauan yang dibangun dapat membantu pengelola kawasan Malioboro memperoleh informasi kepadatan pengunjung secara waktu nyata sebagai dasar pengambilan keputusan pengendalian kerumunan.
    \item Pemantauan berbasis \textit{computer vision} memberikan informasi kepadatan yang lebih objektif dan mencakup lebih banyak titik dibandingkan pengamatan manual oleh operator.
\end{enumerate}

\section{Sistematika Penulisan}

Bab I berisi pendahuluan yang menjelaskan dasar dilakukannya penelitian. Bagian latar 
belakang memaparkan kondisi kepadatan pengunjung di kawasan Malioboro dan keterbatasan 
pemantauan manual yang berjalan saat ini, yang menjadi alasan dibangunnya sistem 
pemantauan otomatis. Dari latar belakang tersebut diturunkan rumusan masalah, tujuan, 
batasan, dan manfaat penelitian. Bab ini ditutup dengan sistematika penulisan yang 
menggambarkan susunan keseluruhan naskah.

Bab II menyajikan tinjauan pustaka dan dasar teori yang melandasi penelitian. Tinjauan 
pustaka membahas penelitian terdahulu yang relevan beserta posisi penelitian ini di 
antaranya. Dasar teori menguraikan konsep deteksi objek, arsitektur model YOLOv8 dan 
YOLO11, metrik evaluasi deteksi, ketidakseimbangan kelas, anotasi semi-otomatis, dan 
tingkat pelayanan (\textit{level of service}) pejalan kaki. Bab ini ditutup dengan 
analisis perbandingan metode yang menjadi dasar pemilihan model deteksi.

Bab III menguraikan metode penelitian secara bertahap. Pembahasan dimulai dari pengumpulan 
data berupa aliran kamera pengawas, dilanjutkan anotasi semi-otomatis dengan koreksi 
manual dan prapemrosesan data untuk membentuk dataset. Bagian berikutnya menjelaskan 
konfigurasi pelatihan model serta metode evaluasi yang digunakan untuk membandingkan 
varian model, serta penyetelan \textit{hyperparameter} pada model terpilih untuk 
mengoptimasi performanya. Bab ini juga memuat perancangan dan implementasi sistem 
pemantauan berbasis web yang memproses video secara waktu nyata.

Bab IV menyajikan hasil penelitian beserta pembahasannya. Bagian awal memaparkan hasil 
pembentukan dataset, dilanjutkan hasil pelatihan dan perbandingan varian model YOLOv8 dan 
YOLO11 berdasarkan metrik yang telah ditetapkan. Analisis performa pada tiap kelas 
dibahas untuk melihat pengaruh ketidakseimbangan kelas terhadap hasil deteksi. Bagian 
selanjutnya menyajikan hasil penyetelan \textit{hyperparameter}
pada model terpilih beserta dampaknya terhadap akurasi dan kecepatan inferensi. Bagian
akhir memaparkan hasil implementasi dan pengujian sistem pemantauan.

Bab V berisi kesimpulan dan saran. Kesimpulan menjawab setiap rumusan masalah berdasarkan 
temuan yang diperoleh pada Bab IV. Saran memuat arah pengembangan yang dapat dilakukan 
pada penelitian selanjutnya sesuai dengan lingkup penelitian ini.